\chapter{地址转换、TLB 与虚拟内存}

\begin{keyidea}
处理器发出的虚拟地址必须经过权限、页表和 TLB 才能成为实际事务。本章用 80386 的历史边界进入现代地址转换路径。
\end{keyidea}

\section{一、80386：32-bit、保护模式与地址转换}

80386 把 x86 推入 32-bit 时代。它扩大了通用寄存器宽度，引入 EAX、EBX、ECX、EDX、ESP、EBP、ESI、EDI 等 32-bit 形式；地址和数据处理能力显著增强，并提供保护模式、分页支持和更完整的内存保护机制。相较 8086，80386 开始承担更复杂的系统级硬件职责。

80386 的一个关键变化是 32-bit 线性地址。程序形成的有效地址先经过分段机制得到线性地址；在启用分页时，线性地址再经过页表转换为物理地址。可以用简化链路表示：

\begin{lstlisting}
// 注释（推进）：按行追踪数据来源、经过的部件和最终写入目标。
// 注释（边界）：只有标出的时钟沿、阶段或异常入口会改变体系结构可见状态。
effective address -> segmentation -> linear address
linear address    -> paging       -> physical address
\end{lstlisting}

这条链路说明，地址不再只是“寄存器加偏移后送到总线”。CPU 内部必须有地址形成、权限检查、页表访问、异常触发等硬件机制。若页表项不存在、权限不允许、访问越界，CPU 产生异常事件，控制流转入规定入口。

保护模式把“谁能访问什么”变成硬件可检查的约定。段描述符、特权级、页表权限等机制，使得不同程序或系统组件不能随意访问全部内存或执行全部指令。这些机制后来成为现代操作系统隔离、虚拟内存和系统调用的硬件基础。CPU 不只执行 ALU 运算，还参与地址翻译、权限判断和异常入口选择。

\begin{longtable}{L{0.24\linewidth}L{0.58\linewidth}}
\toprule
80386 观察点 & 硬件含义 \\
\midrule
32-bit 通用寄存器 & 数据通路和编程模型扩展到 32-bit \\\tablerowrule
32-bit 线性地址 & 地址空间显著扩大，地址形成更复杂 \\\tablerowrule
保护模式 & CPU 硬件检查权限、界限和特权级 \\\tablerowrule
分页机制 & 线性地址可经页表映射到物理地址 \\\tablerowrule
异常机制 & 地址或权限错误进入硬件规定的控制入口 \\
\bottomrule
\end{longtable}

80386 的外部接口同样要按事务理解。32-bit 数据路径并不意味着每次访问都天然是完整 32-bit 字。处理器需要用地址低位和字节使能说明本次访问哪些字节有效；外部总线、存储器阵列和 cache 控制器必须据此选择正确字节、处理未对齐访问或拆分周期。

\begin{longtable}{L{0.22\linewidth}L{0.34\linewidth}L{0.34\linewidth}}
\toprule
80386 对象 & 硬件含义 & 常见调试点 \\
\midrule
地址线与数据线 & 指出事务目标并传输 32-bit 数据片段 & 地址对齐、字节选择和目标译码是否一致 \\\tablerowrule
字节使能 & 表明本周期哪些字节有效 & 字节/半字/字访问不会破坏相邻字节 \\\tablerowrule
总线状态 & 区分读、写、取指、确认等周期语义 & 外部控制器是否按状态产生正确控制线 \\\tablerowrule
分页部件 & 查页目录和页表，形成物理地址 & 缺页、权限和存在位错误是否触发异常 \\\tablerowrule
TLB 思想 & 缓存近期地址转换结果 & 命中降低转换时延，失效和权限改变要刷新 \\
\bottomrule
\end{longtable}

分页可以用一次内存加载来观察。指令先形成有效地址，分段得到线性地址；分页开启时，线性地址被拆成页目录索引、页表索引和页内偏移。若转换结果已在 TLB 中，地址转换可以很快完成；若未命中，就要访问内存中的页表结构。页表项还带有存在位、读写权限、用户/内核权限、访问位和脏位等硬件标志。任何一步失败，这次访问不应返回任意数据，而应进入异常处理路径。

\begin{lstlisting}
// 注释（推进）：逐行追踪发起者、所有权和可见性；请求被接受不等于事务完成。
// 注释（边界）：以采样沿、状态位或 completion 为准，再允许消费者使用结果。
80386-style address walk:
  effective address
    -> segmentation checks
    -> linear address
    -> TLB lookup or page-table walk
    -> permission checks
    -> physical address or exception
\end{lstlisting}

\begin{pdfdiagram}
  \node[hw state, minimum width=1.8cm] (ea) at (0,0.6) {有效地址};
  \node[hw compute, minimum width=1.95cm] (seg) at (2.35,0.6) {分段检查};
  \node[hw state, minimum width=1.95cm] (lin) at (4.8,0.6) {线性地址};
  \node[hw compute, minimum width=1.95cm] (tlb) at (7.25,0.6) {TLB/页表};
  \node[hw state, minimum width=1.95cm] (pa) at (9.7,0.6) {物理地址};
  \node[hw control, minimum width=2.2cm] (ex) at (7.25,-1.15) {异常入口};
  \draw[hw bus] (ea) -- node[above=2pt,hw label,fill=white,inner sep=1pt]{选择子/偏移} (seg);
  \draw[hw bus] (seg) -- node[above=2pt,hw label,fill=white,inner sep=1pt]{界限/权限通过} (lin);
  \draw[hw bus] (lin) -- node[above=2pt,hw label,fill=white,inner sep=1pt]{PDE/PTE} (tlb);
  \draw[hw bus] (tlb) -- node[above=2pt,hw label,fill=white,inner sep=1pt]{允许访问} (pa);
  \draw[hw return] (seg.south) |- node[pos=0.45,left,hw label,fill=white,inner sep=1pt]{段错误} (ex.west);
  \draw[hw return] (tlb.south) -- node[right,hw label,fill=white,inner sep=1pt]{页错误} (ex.north);
  \node[hw label, text width=9.5cm] at (4.85,-2.15) {80386 的一次内存加载在真正访问 cache 或外部总线前，已经可能因为段、页或权限检查失败而进入异常路径。};
\end{pdfdiagram}
\diagramnote{80386 地址转换链路。地址转换和权限检查是访存事务的一部分，不是软件注释。}

\topic{用一个线性地址拆开分页硬件}以线性地址 \code{0x00403ABC} 为例，80386 的 4 KiB 页式分页可把它拆成页目录索引、页表索引和页内偏移。高 10 bit 选择页目录项，中间 10 bit 选择页表项，低 12 bit 是页内偏移。该地址的页目录索引为 1，页表索引为 3，页内偏移为 \code{0xABC}。若 CR3 指向的页目录中第 1 项不存在，访问在真正读数据前就会变成异常；若页表项存在但权限不允许，也不能继续把物理地址送到 cache 或外部总线。

\begin{longtable}{L{0.2\linewidth}L{0.24\linewidth}L{0.46\linewidth}}
\toprule
字段 & 示例值 & 硬件用途 \\
\midrule
页目录索引 & 1 & 选择页目录中的一个 PDE \\\tablerowrule
页表索引 & 3 & 选择对应页表中的一个 PTE \\\tablerowrule
页内偏移 & \code{0xABC} & 与物理页框基址相加形成最终地址 \\\tablerowrule
P/RW/US 等位 & 来自 PDE/PTE & 判断存在性、读写权限和用户/特权访问 \\
\bottomrule
\end{longtable}

\topic{把这次拆页走成一次完整的两级 walk}上面的例子只做了字段拆分，还没有真正读页表。现在把 \code{0x00403ABC} 从头到尾走完，假设 CR3 保存的页目录基址是 \code{0x00005000}（页目录本身 4 KiB 对齐，低 12 bit 恒为 0，不参与拼接）。第一步，用页目录索引 1 算出 PDE 的地址：\code{0x00005000 + 1×4 = 0x00005004}，分页部件在这里做一次 4 字节内存读，取回 PDE。设取回的 PDE 为 \code{0x00007027}，其高 20 bit 给出页表基址 \code{0x00007000}，低 12 bit 是标志位 \code{0x027}（P=1、RW=1、US=1、A=1）。第二步，用页表索引 3 算出 PTE 的地址：\code{0x00007000 + 3×4 = 0x0000700C}，再做一次 4 字节内存读，取回 PTE。设取回的 PTE 为 \code{0x0002B027}，其高 20 bit 给出物理页框基址 \code{0x0002B000}，标志位同样允许本次访问。第三步，把页内偏移 \code{0xABC} 拼上：\code{0x0002B000 + 0xABC = 0x0002BABC}，这才是最终送到 cache 和外部总线的物理地址。

\begin{lstlisting}
CR3             = 0x00005000           (page directory base)
PDE address     = CR3 + 1*4            = 0x00005004
PDE             = 0x00007027           (table base 0x00007000, flags 0x027)
PTE address     = 0x00007000 + 3*4     = 0x0000700C
PTE             = 0x0002B027           (frame base 0x0002B000, flags 0x027)
physical addr   = 0x0002B000 + 0xABC   = 0x0002BABC
\end{lstlisting}

这次 walk 的代价值得单独看一眼：在读到目标数据之前，分页部件已经先做了两次内存读（一次读 PDE，一次读 PTE）。若每次访存都这样走，内存流量和地址时延都会成倍增加。TLB 解决的正是这个问题：近期用过的“线性页号到物理页框”转换缓存在片内，命中时两次页表读全部省掉，只剩目标数据那一次访问。权限位和存在位也要随转换结果一起缓存，否则 TLB 命中路径会绕过检查；反过来，软件修改页表后必须刷新对应的 TLB 项，否则硬件会继续按旧映射放行访问。这也是存储与互连相关章节讨论 cache 时见过的“缓存必然带来失效问题”在地址转换路径上的重演。

\begin{longtable}{L{0.18\linewidth}L{0.34\linewidth}L{0.36\linewidth}}
\toprule
walk 步骤 & 地址算术 & 关键问题 \\
\midrule
读 PDE & \code{0x00005000 + 1×4 = 0x00005004} & PDE 的 P 位是否为 1，页表是否存在 \\\tablerowrule
读 PTE & \code{0x00007000 + 3×4 = 0x0000700C} & PTE 的 P/RW/US 是否允许本次访问 \\\tablerowrule
拼物理地址 & \code{0x0002B000 + 0xABC = 0x0002BABC} & 页框基址加偏移，不再查表 \\\tablerowrule
TLB 命中 & 不访问页表，直接得到页框 & 两次页表读全部省掉，只读目标数据 \\
\bottomrule
\end{longtable}

\begin{pdfdiagram}
  \node[hw state, minimum width=3.0cm] (va) at (0,0) {线性地址 \code{0x00403ABC}};
  \node[hw compute, minimum width=1.2cm] (tlb) at (4.3,0) {TLB};
  \node[hw state, minimum width=2.8cm] (cr3) at (0,-1.5) {CR3 = \code{0x00005000}};
  \node[hw memory, minimum width=2.8cm] (pde) at (0,-3.0) {PDE = \code{0x00007027}};
  \node[hw memory, minimum width=2.8cm] (pte) at (0,-4.5) {PTE = \code{0x0002B027}};
  \node[hw memory, minimum width=2.8cm] (frm) at (0,-6.0) {页框基址 \code{0x0002B000}};
  \node[hw state, minimum width=2.8cm] (pa) at (4.3,-6.0) {物理地址 \code{0x0002BABC}};
  \draw[hw bus] (va) -- node[above=1pt,hw label]{页号 \code{0x00403}} (tlb);
  \draw[hw return] (tlb.south) -- node[right=1pt,hw label]{命中：直接得页框} (pa.north);
  \draw[hw bus] (va.south) -- node[midway,left=2pt,hw label]{未命中：开始 walk} (cr3.north);
  \draw[hw bus] (cr3) -- node[midway,left=2pt,hw label]{读 \code{0x5000+1×4=0x5004}} (pde);
  \draw[hw bus] (pde) -- node[midway,left=2pt,hw label]{读 \code{0x7000+3×4=0x700C}} (pte);
  \draw[hw bus] (pte) -- node[midway,left=2pt,hw label]{高 20 bit 给页框} (frm);
  \draw[hw bus] (frm) -- node[above=1pt,hw label,fill=white,inner sep=1pt]{+ 偏移 \code{0xABC}} (pa);
  \node[hw label] at (4.3,-6.6) {送 cache 与外部总线};
\end{pdfdiagram}
\diagramnote{80386 两级页表 walk 的指针链（线性地址 \code{0x00403ABC}：页目录索引 1、页表索引 3、页内偏移 \code{0xABC}）。CR3 给出页目录基址 \code{0x5000}，读 \code{0x5004} 取回 PDE=\code{0x00007027}；按页表基址 \code{0x7000} 读 \code{0x700C} 取回 PTE=\code{0x0002B027}；页框 \code{0x0002B000} 加偏移得物理地址 \code{0x0002BABC}。右侧 TLB 命中时整条链被旁路，两次页表读全部省掉。}

\topic{控制寄存器和描述符表是系统约定的入口}80386 的保护模式不能只理解为“地址变大”。处理器必须知道段描述符在哪里、异常入口在哪里、页目录在哪里、当前是否启用保护和分页。这些信息由控制寄存器、GDTR/IDTR 等状态保存。它们是 CPU 执行内存访问、异常和特权检查时自动使用的硬件约定。

\begin{longtable}{L{0.18\linewidth}L{0.32\linewidth}L{0.4\linewidth}}
\toprule
对象 & 保存内容 & 硬件使用场景 \\
\midrule
CR0 & 保护模式、分页等控制位的一部分 & 决定地址和权限检查的总体模式 \\\tablerowrule
CR3 & 页目录基址 & TLB 未命中或刷新后查页表 \\\tablerowrule
GDTR & 全局描述符表位置和界限 & 装载段选择子、检查段属性 \\\tablerowrule
IDTR & 中断/异常描述符表位置和界限 & 异常、中断、陷入入口选择 \\\tablerowrule
EFLAGS & 条件码、方向、中断允许等状态 & 分支、算术结果和中断接收边界 \\
\bottomrule
\end{longtable}

\topic{从实模式进入保护模式是一段硬件约定切换}80386 复位后并不会自动处于完整分页保护环境。系统固件或早期启动代码必须先准备描述符表，再装载 GDTR，设置 CR0 中的保护模式相关位，并通过跳转刷新取指路径和段状态。分页还需要准备页目录、页表，把 CR3 指向页目录，再打开分页控制位。每一步都是硬件状态切换，顺序错误会导致取指、数据访问或异常入口无法解释。

\begin{longtable}{L{0.18\linewidth}L{0.4\linewidth}L{0.32\linewidth}}
\toprule
阶段 & 硬件动作 & 错误后果 \\
\midrule
准备 GDT/IDT & 在内存中放置段描述符和异常入口描述符 & 段装载或异常入口失败 \\\tablerowrule
装载 GDTR/IDTR & CPU 记录描述符表基址和界限 & 后续选择子无法被正确解释 \\\tablerowrule
设置 CR0 & 开启保护模式相关控制 & 旧的取指和段状态需要刷新 \\\tablerowrule
远跳转/重装段 & 让 CS/DS/SS 等进入新描述符语义 & 继续按旧语义执行会失控 \\\tablerowrule
准备页表并设置 CR3 & 给分页硬件页目录入口 & TLB 未命中时无法 walk \\\tablerowrule
开启分页 & 线性地址开始经页表转换 & 映射缺失会立即异常 \\
\bottomrule
\end{longtable}

\begin{pdfdiagram}
  \node[hw state] (real) at (0,0) {实模式};
  \node[hw compute, align=center, minimum width=2.5cm] (gdtp) at (3.15,0) {① 内存中准备\\GDT 与 IDT};
  \node[hw compute, align=center, minimum width=2.5cm] (lgdt) at (6.5,0) {② LGDT / LIDT\\装载 GDTR / IDTR};
  \node[hw compute, align=center, minimum width=2.5cm] (pe) at (6.5,-1.9) {③ 置 CR0 的\\PE 位};
  \node[hw compute, align=center, minimum width=2.5cm] (jmp) at (3.15,-1.9) {④ 远跳转刷新\\CS 与取指路径};
  \node[hw state] (prot) at (0,-1.9) {保护模式};
  \draw[hw bus] (real) -- (gdtp);
  \draw[hw bus] (gdtp) -- (lgdt);
  \draw[hw bus] (lgdt.south) -- (pe.north);
  \draw[hw bus] (pe) -- (jmp);
  \draw[hw bus] (jmp) -- (prot);
  \node[hw label, text width=9.6cm] at (3.4,-3.2) {顺序不能颠倒：GDT 没准备好，LGDT 装载的就是无效表基址；PE 置位改变了段的解释方式，必须靠一次远跳转让 CS 和取指路径进入新语义。分页在此之后：再准备页表、把 CR3 指向页目录、开启分页位。};
\end{pdfdiagram}
\diagramnote{从实模式进入保护模式的硬件步骤链（80386）。准备描述符表、装载 GDTR/IDTR、置 CR0 的 PE 位、远跳转刷新取指路径与段状态，四步都是硬件状态切换，顺序错误会让取指、数据访问或异常入口无法解释；若还要分页，需在此之后再准备页目录和页表、装载 CR3、开启分页控制位。}

\topic{描述符和异常错误码把失败原因带进硬件入口}保护模式中的段描述符不只是“基址和长度”。它还包含类型、存在位、特权级、可读写/可执行等属性。分页项也不只是物理地址，还包含存在、读写、用户/特权、访问和脏等位。异常发生时，CPU 需要把足够信息交给异常入口，使系统能够判断是不存在、权限错误、越界，还是其他硬件条件失败。

\begin{longtable}{L{0.2\linewidth}L{0.34\linewidth}L{0.36\linewidth}}
\toprule
对象 & 关键字段 & 硬件问题 \\
\midrule
段描述符 & base、limit、type、DPL、present & 这个选择子是否存在、能否以当前特权使用 \\\tablerowrule
页表项 & page frame、present、RW、US、accessed、dirty & 这个线性页是否存在、能否按本次方式访问 \\\tablerowrule
异常入口 & selector、offset、gate type、DPL & 失败后 CPU 应进入哪段处理代码 \\\tablerowrule
错误码 & 选择子或页错误原因位 & 处理程序如何定位失败对象和原因 \\
\bottomrule
\end{longtable}

一次页错误不是“内存读失败”这么简单。CPU 已经形成了线性地址，检查 PDE/PTE 或权限时发现不能继续，于是放弃普通写回路径，保存足够的异常信息，并从 IDT 中选择页错误入口。若处理器允许一部分状态已经提交、一部分状态未提交，系统就难以恢复。因此异常路径必须与指令提交边界配合：错误指令不应像成功指令一样修改目标寄存器或设备状态。

\topic{页错误诊断要同时看地址、访问类型和权限位}80386 把页错误从“总线读不到数据”提升为可诊断的硬件事件。CPU 会保存导致故障的线性地址，并通过异常错误码给出若干原因位。学习阶段不必记住所有后续扩展，但必须掌握三个基本问题：页是否存在，访问是读还是写，访问来自用户态还是特权态。只有把这三项与 PDE/PTE 中的 present、RW、US 等位对照，才能判断是缺页、写只读页、用户态访问内核页，还是页表本身没有建立。

\begin{longtable}{L{0.2\linewidth}L{0.34\linewidth}L{0.36\linewidth}}
\toprule
观察对象 & 含义 & 排查方向 \\
\midrule
故障线性地址 & 哪个线性地址无法完成转换或权限检查 & 页目录索引、页表索引和页内偏移是否落在预期区域 \\\tablerowrule
P 位 & 故障由不存在页还是保护违规引起 & 缺页处理、页表初始化或权限设置 \\\tablerowrule
W/R 位 & 本次访问是写还是读导致故障 & 写只读页、写保护代码段、copy-on-write 类机制 \\\tablerowrule
U/S 位 & 本次访问来自用户级还是特权级 & 用户程序越权访问内核映射或段/页权限不一致 \\\tablerowrule
IDT 入口 & 页错误应跳到哪个处理程序 & 异常入口描述符、栈切换和处理代码是否有效 \\
\bottomrule
\end{longtable}

这套诊断方法可以反过来指导最小保护模式实验。先建立一段恒等映射，让取指、栈和数据访问都可成功；再故意取消某个页表项的 present 位，确认访问该页会进入页错误入口；随后把某页设置为只读，执行写访问，确认错误码能区分权限失败和不存在页。若异常处理程序本身所在页没有映射，CPU 可能在处理一个错误时再次出错，最终表现为无法恢复的故障。因此保护模式 bring-up 必须先保证 GDT、IDT、栈、异常代码和页表所在内存可访问。

\topic{一次页错误可以把三个原因位逐项读完}看一个完整案例。某次内存访问触发页错误，CPU 把故障线性地址保存在 CR2：\code{CR2 = 0x00DEAD00}；异常入口处的错误码为 \code{0x00000003}。先拆 CR2：\code{0x00DEAD00 >> 12 = 0x00DEA}，说明故障发生在线性页 \code{0x00DEA} 内偏移 \code{0xD00} 处。再拆错误码低 3 bit：bit0（P）=1，bit1（W/R）=1，bit2（U/S）=0。逐项解释。P=1 表示保护违规：PDE/PTE 链存在，CPU 是在权限检查上拦下了这次访问。W/R=1 表示触发故障的是一次写访问。U/S=0 表示访问发生时 CPU 处于特权级，即内核态代码。三项合起来，结论只有一句：\hwkey{内核态写只读页}。典型场景是内核错误地写了一个按只读映射的页，例如写保护的常量区或被标记为 copy-on-write 的页。这里有一个历史细节值得知道：80386 上特权级访问不受 RW 位约束，80486 起 CR0 增加 WP 位，置位后特权态写只读页同样触发页错误；不置 WP 时这类错误可能悄悄写成数据损坏，置 WP 之后它才变成可诊断的页错误。

同样的读法可以做两个练习。第一，错误码 \code{0x00000004}：P=0、W/R=0、U/S=1，即用户态的一次读访问命中了不存在的页。此时看 CR2：若它是个很小的地址（例如 \code{0x00000018}），最常见的解释是用户程序解引用了空指针或野指针；若它落在合法但尚未分配的堆、栈区域，则是典型的按需分配缺页，处理程序应分配页框、建立映射后重新执行该指令。第二，错误码 \code{0x00000007}：P=1、W/R=1、U/S=1，即用户态写一个已存在但只读的用户页，典型解释是程序试图改写常量数据或共享库的只读映射；若系统正在使用 copy-on-write，这反而是预期中的缺页，处理程序复制页框后放行即可。判断顺序永远相同：先看 P 区分“不存在”与“保护违规”，再看 W/R 确定访问类型，然后用 U/S 确定访问来源，最后拿 CR2 对照页表定位具体对象。

\begin{pdfdiagram}
  \node[hw state, minimum width=5cm] (root) at (0,0) {页错误发生\\读 CR2 与错误码低 3 位：P、W/R、U/S};
  \node[hw control, minimum width=4.2cm] (q1) at (0,-1.9) {① P 位：页存在吗？};
  \node[hw control, minimum width=4.2cm] (q2) at (0,-3.7) {② W/R 位：写访问吗？};
  \node[hw control, minimum width=4.2cm] (q3) at (0,-5.5) {③ U/S 位：访问来自哪一级？};
  \node[hw memory, text width=4.6cm] (c1) at (6.9,-1.9) {缺页（P=0）\\对照 CR2：很小的地址多为空指针；未分配区则分配页框后重试};
  \node[hw memory, text width=4.6cm] (c4) at (6.9,-3.7) {读侧违规（P=1, W/R=0）\\多为用户态读特权页：U/S=1 访问 U/S=0 页};
  \node[hw memory, text width=4.6cm] (c2) at (6.9,-5.5) {内核态写只读页（U/S=0）\\写保护的常量区或被标记 copy-on-write 的页};
  \node[hw memory, text width=4.6cm] (c3) at (6.9,-7.3) {用户态写只读页（U/S=1）\\改写常量数据；copy-on-write 则复制页框后放行};
  \draw[hw return] (root.south) -- (q1.north);
  \draw[hw return] (q1.east) -- node[above=1pt,hw label,fill=white,inner sep=1pt]{P=0} (c1.west);
  \draw[hw return] (q1.south) -- node[right,hw label,fill=white,inner sep=1pt]{P=1：保护违规} (q2.north);
  \draw[hw return] (q2.east) -- node[above=1pt,hw label,fill=white,inner sep=1pt]{W/R=0 读} (c4.west);
  \draw[hw return] (q2.south) -- node[right,hw label,fill=white,inner sep=1pt]{W/R=1 写} (q3.north);
  \draw[hw return] (q3.east) -- node[above=1pt,hw label,fill=white,inner sep=1pt]{U/S=0 内核态} (c2.west);
  \draw[hw return] (q3.south) -- (0,-7.3) -- node[pos=0.7,above=1pt,hw label,fill=white,inner sep=1pt]{U/S=1 用户态} (c3.west);
\end{pdfdiagram}
\diagramnote{页错误诊断决策树：先看 P 位区分缺页与保护违规，缺页时拿 CR2 对照预期区域；保护违规再看 W/R 位确定读或写，最后用 U/S 位确定访问来源。左列是固定的判断次序，右列是四种结论与各自的处理方向。}

\begin{longtable}{L{0.16\linewidth}L{0.30\linewidth}L{0.42\linewidth}}
\toprule
错误码 & 三个原因位 & 判断与排查方向 \\
\midrule
\code{0x00000003} & P=1、W/R=1、U/S=0 & 特权态写只读页；检查该页映射属性和内核写路径 \\\tablerowrule
\code{0x00000004} & P=0、W/R=0、U/S=1 & 用户态读不存在页；空指针、野指针或按需分配缺页 \\\tablerowrule
\code{0x00000007} & P=1、W/R=1、U/S=1 & 用户态写只读用户页；常量区误写或 copy-on-write \\
\bottomrule
\end{longtable}

\topic{80386 总线的字节使能与 8086 一脉相承}80386 对外有 32-bit 数据总线。地址线通常给出 32-bit 对齐字的高位地址，字节使能信号指出本次传输中哪些 8-bit 字节有效。这样，一个字节写不会破坏同一 32-bit 数据片段中的其他字节；一个未对齐的 32-bit 写可能需要拆成多个传输，或由硬件和总线控制器处理额外周期。

\begin{longtable}{L{0.2\linewidth}L{0.28\linewidth}L{0.42\linewidth}}
\toprule
访问类型 & 有效字节车道 & 关键问题 \\
\midrule
字节访问 & 一个 BE 低有效 & 相邻字节不能被写改 \\\tablerowrule
半字访问 & 两个相邻 BE 低有效 & 地址最低位和字节序必须一致 \\\tablerowrule
对齐字访问 & 四个 BE 低有效 & 一个总线数据片段可完成传输 \\\tablerowrule
未对齐字访问 & 跨越两个对齐片段 & 需要拆分周期或触发规定异常 \\
\bottomrule
\end{longtable}

80486 把片上 cache、流水线和浮点方向的集成推到更显著的位置，使常用指令和数据更容易走近路径；Pentium 进一步展示双流水线、分支预测和更复杂控制对吞吐率的影响。它们复制、切分和调度执行路径，让常见工作更少等待远处存储器和长控制路径。

\begin{longtable}{L{0.2\linewidth}L{0.32\linewidth}L{0.38\linewidth}}
\toprule
阶段 & 新增硬件重点 & 与前文概念的连接 \\
\midrule
80386 & 32-bit 地址、分页、保护和异常 & 地址路径和权限判断进入 CPU 内部关键结构 \\\tablerowrule
80486 & 片上 cache、流水线、浮点集成 & 常见数据靠近核心，长组合路径被边界切分 \\\tablerowrule
Pentium & 双流水线、预测和更复杂发射控制 & 多条候选执行路径需要统一调度和恢复 \\\tablerowrule
现代核心 & 乱序、重命名、TLB、多级 cache、多核互连 & 快慢路径、状态提交和一致性成为核心问题 \\
\bottomrule
\end{longtable}

\section{二、多级页表与 TLB 的现代路径}

虚拟地址由虚拟页号 VPN 和页内偏移组成。页大小为 4 KiB 时偏移占 12 bit；地址转换只替换 VPN，页内偏移原样进入物理地址。多级页表把 VPN 再切成若干索引，只为实际使用的地址区间分配下级表，避免一张巨大平面表常驻内存。

以四级页表为例，一次 TLB miss 最坏要读取四个页表项，再做真正的数据访问。硬件 page-table walker 发出的这些读请求也经过 Cache；因此页表局部性、页表项权限和普通数据 Cache 共同决定 miss 代价。大页减少页表层数和 TLB 压力，却增大内部碎片和迁移成本。

\begin{longtable}{L{0.2\linewidth}L{0.34\linewidth}L{0.36\linewidth}}
\toprule
事件 & 谁处理 & 结果 \\
\midrule
TLB hit & MMU 硬件 & 立即得到页框号并检查权限 \\\tablerowrule
TLB miss、页表项有效 & page-table walker & 读取页表并回填 TLB \\\tablerowrule
页表项无效但可调页 & 操作系统缺页处理 & 分配或从存储读入页面后重试 \\\tablerowrule
权限失败或地址非法 & 异常处理程序 & 向进程报告错误或终止 \\
\bottomrule
\end{longtable}

\section{三、上下文切换、ASID 与失效协议}

切换进程时，旧 TLB 项不能被新地址空间误用。最简单的办法是清空 TLB；带 ASID/PCID 的实现把地址空间标识一同参与匹配，使多个进程的翻译可以共存。操作系统修改已缓存的页表项后，还必须向相关核心发出 TLB shootdown，并等待它们确认失效，之后才能回收旧页框。

设备侧 IOMMU 使用相同思想把设备虚拟地址翻译成物理页框，并按设备或进程检查权限。它解决的是 DMA 隔离与地址重映射，不会自动替代 CPU Cache 一致性和描述符可见性屏障。

\section{四、逐位走完一次四级页表遍历}

考虑一种常见的 48-bit 虚拟地址、4 KiB 页面、每级页表含 512 项的四级结构。页内偏移占 12 bit，每级索引占 9 bit，虚拟地址可写成 $[L4:9][L3:9][L2:9][L1:9][offset:12]$。根页表的物理地址保存在页表基址寄存器中；walker 用 L4 索引读取第一级页表项，取出下一级页框，再依次使用 L3、L2、L1 索引。末级页表项给出物理页框号，与原 offset 拼接成物理地址。

每次读取页表项都必须检查 present、读写、用户/特权、可执行和保留位等约束。权限通常沿层级收紧：上层禁止用户访问时，下层不能重新放开。若中间层页表项声明大页，遍历提前结束，由更低的部分 VPN 直接充当大页内偏移。由此可见，大页既减少 TLB 项需求，也减少 walker 的内存访问次数。

\begin{pdfdiagram}[
  node distance=7mm and 9mm,
  vmbox/.style={draw, rounded corners, align=center, minimum width=25mm, minimum height=9mm, fill=blue!5},
  vmdecision/.style={draw, diamond, aspect=2.0, align=center, fill=yellow!8},
  >=Stealth]
\node[vmbox] (req) {虚拟地址访问};
\node[vmdecision, right=of req] (tlb) {TLB 命中？};
\node[vmbox, right=of tlb, yshift=10mm] (access) {权限检查\\访问 Cache};
\node[vmbox, right=of tlb, yshift=-10mm] (walk) {逐级读取 PTE};
\node[vmdecision, right=of walk] (valid) {映射有效？};
\node[vmbox, right=of valid, yshift=10mm] (fill) {回填 TLB};
\node[vmbox, right=of valid, yshift=-10mm] (fault) {精确缺页异常\\操作系统处理};
\draw[->] (req) -- (tlb);
\draw[->] (tlb) -- node[above, font=\small]{是} (access);
\draw[->] (tlb) -- node[below, font=\small]{否} (walk);
\draw[->] (walk) -- (valid);
\draw[->] (valid) -- node[above, font=\small]{是} (fill);
\draw[->] (fill) |- (access);
\draw[->] (valid) -- node[right, font=\small]{否} (fault);
\draw[->] (fault) |- node[pos=0.25, below, font=\small]{修复映射后重试} (req);
\end{pdfdiagram}
\diagramnote{TLB miss 后的地址转换与重试路径。硬件遍历成功时回填并继续访问；映射缺失时形成精确异常，操作系统修复映射后回到原指令。}

walker 并非每次都从 DRAM 读取四次。页表项参与普通 Cache 层次，一些实现还设置 page-walk cache，缓存上层页表项或部分遍历结果。若多个 miss 共享上层索引，后续遍历可复用状态；但 walker 槽位和内存请求队列有限，密集 TLB miss 仍会形成排队。

\section{五、TLB 层次、权限与访问状态}

处理器常为指令和数据各设小型一级 TLB，再用容量更大的二级统一 TLB 降低 page walk 频率。TLB 项至少保存 VPN、页框号、页大小、ASID/PCID、权限和内存类型。匹配时既要比较地址，也要比较地址空间标识；大页与小页项共存时，还要按各自页大小屏蔽不同数量的低位。

页表项的 accessed 和 dirty 状态帮助操作系统判断页面活跃度与回写需要。有的体系结构由硬件置位，有的由软件通过首次访问异常模拟。更新这些位本身是共享内存写，walker、Cache 和操作系统清位之间必须遵守体系结构规定，避免丢失一次并发访问。

内存类型同样是翻译结果的一部分。普通可缓存内存允许推测和合并，而设备内存往往限制重排、推测或访问宽度。把 MMIO 页错误地映射为普通 Cache 内存，不只是性能问题，还可能使寄存器读写次数和顺序偏离设备协议。

\section{六、VIPT Cache 的别名边界}

L1 常用虚拟 index 与物理 tag 并行 TLB 查询。要避免同一物理地址因不同虚拟别名落入不同集合，index 与 line offset 最好完全位于页内偏移。例如 4 KiB 页面、64 B line 只有 6 bit 可用于 index，因此每路最多容纳 $2^6\times64=4$ KiB；8 路 Cache 可做到 32 KiB 而不越过这一自然边界。

若容量或索引超出页内位，硬件需要额外的别名检测、页着色或其他约束。这里的关键不是“虚拟地址能否访问同一页面”，而是 Cache 在物理 tag 尚未返回时用哪些位选择了集合。指令自修改、同一页多重映射和 DMA 所有权交接都会使别名问题更加明显。

\section{七、缺页异常、写时复制与精确重试}

page fault 是地址转换的受控分支，不等同于程序崩溃。需求零页可分配新物理页并清零；文件映射缺页可从页缓存或存储装入；换出页需要发起 I/O；写时复制则在首次写入共享只读页时分配新页、复制内容并修改当前进程映射。只有地址非法或权限无法修复时，操作系统才把错误报告给进程。

为了正确重试，处理器必须保留发生异常的指令边界，使该指令没有留下对体系结构可见的半完成状态。操作系统修复页表、执行必要的 TLB 失效后，从原指令重新执行。若缺页发生在一次跨页访问中，实现还要确保前半段不会被重复提交两次。

缺页处理中的 I/O 可以让当前线程睡眠，调度器运行其他任务。页面到达后，内核更新 PTE，再将线程置为可运行。因而一次 load 的墙钟时间可能包含存储设备延迟，但执行单元并不会一直占住流水线等待。

\section{八、TLB shootdown 的完整时间线}

修改映射的核心不能在写完 PTE 后立即回收旧页框。它先发布新 PTE 或清除旧 PTE，再向可能缓存该地址空间翻译的核心发送处理器间中断；远端核心执行本地 TLB 失效并确认；发起者收齐确认后，才可把旧页框交给其他用途。否则远端核心可能用陈旧翻译访问已重新分配的数据。

ASID/PCID 减少上下文切换时的全量清空，却不能消除同一地址空间页表修改所需的定向失效。为避免标识复用造成旧项误命中，操作系统还要管理 generation：标识绕回或重新分配时，对相关旧项作全局清理。批量 unmap 可合并失效范围和中断次数，但会延迟页框回收。

\section{九、嵌套翻译与 IOMMU}

虚拟机中，客户虚拟地址先经客户页表成为客户物理地址，再经管理程序维护的第二阶段页表成为主机物理地址。一次普通 page walk 中的每个客户页表地址本身还要经过第二阶段翻译，最坏访存次数会成倍增加，因此硬件会缓存组合翻译和中间结果。故障报告必须区分第一阶段权限错误与第二阶段映射缺失，交给不同软件层处理。

IOMMU 也可提供多级翻译、权限和设备标识隔离。设备提交的地址先由 IOMMU 转换，缺失可能导致请求阻塞、页请求消息或 DMA 错误，取决于平台能力。即使 CPU 与设备共享虚拟地址，TLB 失效、设备在途请求、Cache 一致性和页框回收仍要按明确次序闭合。

\begin{longtable}{L{0.2\linewidth}L{0.32\linewidth}L{0.38\linewidth}}
\toprule
故障位置 & 典型含义 & 主要处理者 \\
\midrule
CPU 第一阶段 & 进程页未映射或权限不符 & 客户/主机操作系统 \\\tablerowrule
CPU 第二阶段 & 虚拟机物理页未映射 & 管理程序 \\\tablerowrule
IOMMU 翻译 & 设备地址越界、权限错误或项缺失 & IOMMU 与设备驱动 \\\tablerowrule
内存访问阶段 & 翻译成功后出现总线或 ECC 错误 & 平台 RAS 路径 \\
\bottomrule
\end{longtable}

\section*{章末练习}

\begin{chapterexercise}{1}
地址分解
\exercisecheck{题目}{将一个 48-bit 虚拟地址按四级 4 KiB 页表拆分，说明每段位数及每级 512 项的来源。}
\end{chapterexercise}

\begin{chapterexercise}{2}
性能分析
\exercisecheck{题目}{在不命中任何 TLB 或页表缓存、页表项均命中数据 Cache 与均落到 DRAM 两种情况下，比较一次 load 需要的存储访问。}
\end{chapterexercise}

\begin{chapterexercise}{3}
Cache 联动
\exercisecheck{题目}{4 KiB 页面、64 B line 的 VIPT Cache 为什么每路容量自然受 4 KiB 约束？8 路时可得到多大容量？}
\end{chapterexercise}

\begin{chapterexercise}{4}
异常路径
\exercisecheck{题目}{按顺序描述写时复制 fault 从指令触发到安全重试的状态变化。}
\end{chapterexercise}

\begin{chapterexercise}{5}
并发正确性
\exercisecheck{题目}{某核心撤销映射后立即复用旧页框，为什么即使已写回 PTE 也不安全？给出正确的 shootdown 顺序。}
\end{chapterexercise}

\begin{chapterexercise}{6}
设备隔离
\exercisecheck{题目}{IOMMU 翻译成功是否意味着 DMA 数据一定已与 CPU Cache 保持一致？说明理由。}
\end{chapterexercise}

\section*{参考解答与要点}

\begin{chapteranswer}{1}
地址分解
\answercheck{解答要点}{offset 为 12 bit，其余 36 bit 分成四段各 9 bit；每项 8 B 时，一页可放 $4096/8=512=2^9$ 项。}
\end{chapteranswer}

\begin{chapteranswer}{2}
性能分析
\answercheck{解答要点}{页表项全命中 Cache 时，逻辑上仍有四级遍历和最终数据访问，但下级流量较小；全部落到 DRAM 时，最坏是四次页表读取加一次数据读取，且存在串行依赖。}
\end{chapteranswer}

\begin{chapteranswer}{3}
Cache 联动
\answercheck{解答要点}{offset 6 bit，页内剩余 6 bit 可作 index，共 64 个集合；每路容量为 $64\times64$ B，即 4 KiB，8 路总容量 32 KiB。}
\end{chapteranswer}

\begin{chapteranswer}{4}
异常路径
\answercheck{解答要点}{写权限检查失败并形成精确异常；内核确认 COW，分配和复制新页，更新当前进程 PTE，失效旧 TLB 项，然后返回原指令重试。}
\end{chapteranswer}

\begin{chapteranswer}{5}
并发正确性
\answercheck{解答要点}{远端 TLB 可能仍缓存旧翻译。应先修改 PTE，向相关核心发送失效请求，等待它们执行并确认，最后回收或复用页框。}
\end{chapteranswer}

\begin{chapteranswer}{6}
设备隔离
\answercheck{解答要点}{不能。IOMMU 处理地址与权限；非一致性 DMA 仍需 Cache 维护和屏障，并要等待设备在途事务结束后才能回收页。}
\end{chapteranswer}
